\documentclass[11pt]{article}

% ---------------------------------------------------------------------------
% AEGIS-AT --- Attribution Integrity Benchmark
% Self-contained, arXiv-ready LaTeX source.
% Compiles with pdflatex (run twice for cross-references / hyperref).
% Dependencies are all standard TeX Live / MiKTeX packages.
% ---------------------------------------------------------------------------

\usepackage[utf8]{inputenc}
\usepackage[T1]{fontenc}
\usepackage[margin=1in]{geometry}
\usepackage{amsmath,amssymb}
\usepackage{booktabs}
\usepackage{xcolor}
\usepackage{microtype}
\usepackage{caption}
\usepackage{enumitem}
\usepackage{listings}
\usepackage{lmodern}
\usepackage[hidelinks]{hyperref}
\hypersetup{
  colorlinks=true,
  linkcolor=black,
  citecolor=black,
  urlcolor=blue!55!black,
  pdftitle={AEGIS-AT: Measuring Attribution Integrity Under Sibling Impersonation in Multi-Agent Delegation},
  pdfauthor={Rahul Singh}
}

% Code listing style ---------------------------------------------------------
\definecolor{codebg}{gray}{0.96}
\definecolor{codecomment}{rgb}{0.30,0.45,0.30}
\definecolor{codekw}{rgb}{0.20,0.20,0.60}
\lstdefinestyle{aegis}{
  basicstyle=\ttfamily\footnotesize,
  backgroundcolor=\color{codebg},
  commentstyle=\color{codecomment}\itshape,
  keywordstyle=\color{codekw}\bfseries,
  numbers=none,
  frame=single,
  framesep=4pt,
  rulecolor=\color{gray!40},
  breaklines=true,
  showstringspaces=false,
  captionpos=b,
  xleftmargin=4pt,
  xrightmargin=4pt
}
\lstset{style=aegis}

% Convenience macros ---------------------------------------------------------
\newcommand{\ais}{\textsc{ais}}
\newcommand{\code}[1]{\texttt{#1}}
\newcommand{\actsub}{\code{act.sub}}
\newcommand{\siem}{\code{siem\_action}}

\title{\textbf{AEGIS-AT: Measuring Attribution Integrity Under\\
Sibling Impersonation in Multi-Agent Delegation}\\[4pt]
\large A Reproducible Adversarial Benchmark Showing That RFC 8693 Delegation\\
Can Make Audit Attribution \emph{Worse}, Not Better}

\author{Rahul Singh\thanks{Independent researcher, Jersey City, NJ.
Correspondence: \href{mailto:rahul.rs1397@gmail.com}{rahul.rs1397@gmail.com}
~$\cdot$~ \href{https://github.com/rahulsingh1397}{github.com/rahulsingh1397}.
Code and documentation are dual-licensed: Apache-2.0 (code) and CC~BY~4.0
(documentation).}}

\date{June 2026}

\begin{document}
\maketitle

\begin{abstract}
\noindent
Standards bodies (NIST NCCoE, OpenID Foundation) are actively recommending
cryptographic delegation --- principally OAuth~2.0 Token Exchange (RFC~8693) ---
as the primitive that will give multi-agent AI systems \emph{non-repudiation}:
the property that every consequential action can be traced to a responsible
party. We present \textbf{AEGIS-AT}, a controlled, reproducible, adversarial
benchmark that tests this assumption directly, and find that it fails in a
specific and structural way. AEGIS-AT implements a minimal Security Operations
Center (SOC) triage pipeline --- two sibling agents sharing one scope-gated tool
--- and measures an \textbf{Attribution Integrity Score} (\ais{}): the fraction
of adversarial actions whose claimed \code{(actor, scope, principal\_chain)}
exactly matches an independently observed ground truth. We score \ais{} across
four progressive defense baselines that are \emph{configuration flags over a
single codebase}, so the four scores are genuinely comparable rather than
apples-to-oranges. The result is a \textbf{non-monotonic curve}: attribution is
perfect ($\ais{}=1.0$) under simple per-agent identity, then \emph{regresses to
zero} ($\ais{}=0.0$) the moment RFC~8693 delegation is added, and tamper-evident
logging does not recover it. The cause is not a bug: RFC~8693's ``current
actor'' (\actsub{}) names the agent that \emph{requested} delegated authority,
but in a multi-agent hand-off the agent that \emph{executes} the action can
differ --- and the standard provides no field that records the executor. An
attacker who can shape the content of a single security alert can therefore
cause a high-consequence action to be taken and attributed to the wrong,
lower-privilege sibling, while every cryptographic check passes. The finding is
\emph{categorical} (it succeeds by construction, not with some probability) and
is verified deterministic: each baseline yields byte-identical records across
repeated runs. All \ais{} values are pre-registered in the threat model before
the attack code was written. We argue the gap is structural rather than
topology-specific, name the standardized layer hypothesized to close it
(sender-constrained tokens, RFC~9449 / RFC~8705), and release the full pipeline
(59 tests; reproduces in seconds).
\end{abstract}

\noindent\textbf{Keywords:} multi-agent systems, AI agent security, attribution,
non-repudiation, OAuth 2.0 Token Exchange, RFC 8693, confused deputy, delegation,
audit logging, benchmark.

\tableofcontents

% ===========================================================================
\section{Introduction}
% ===========================================================================

Multi-agent AI systems increasingly take consequential, irreversible actions:
quarantining hosts, blocking accounts, modifying infrastructure. When such an
action is taken, an audit log is expected to answer one question: \emph{which
agent did this?} The dominant answer being standardized is cryptographic
delegation. A human principal authenticates once; an orchestrator mints scoped
delegation tokens for each agent; the token carries a verifiable chain of
``who acts on whose behalf.'' OAuth~2.0 Token Exchange (RFC~8693) is the
mechanism, and NIST's NCCoE and the OpenID Foundation have both, in early 2026,
named auditing and non-repudiation of AI agent actions as an open problem to be
solved with exactly these primitives~\cite{nist2026,openid2026}.

This paper asks whether that assumption survives contact with a realistic
adversary, and answers: \emph{not in the multi-agent re-delegation case}. We
build the smallest system in which the failure can occur --- two sibling agents
(\code{Agent-Enrich}, read-only; \code{Agent-Contain}, write-capable) sharing one
scope-gated tool --- and we measure attribution correctness under a
sibling-impersonation attack across four defense configurations.

\paragraph{The finding.}
Attribution is \emph{perfect} under simple per-agent identity (Baseline~2,
$\ais{}=1.0$). Adding RFC~8693 delegation (Baseline~3) drives it to
$\ais{}=0.0$, and tamper-evident logging (Baseline~4) leaves it there. The curve
is non-monotonic: the two primitives most emphasized for agent non-repudiation
--- signed delegation chains and tamper-evident logs --- do not close the gap,
and signed delegation \emph{actively opens} it relative to plain per-agent
identity. Table~\ref{tab:curve} states the result.

\begin{table}[h]
\centering
\caption{The measured Attribution Integrity Score across four baselines. Each
baseline is a configuration flag over one codebase; only the credential differs.
Measured values reproduce the pre-registered prediction exactly.}
\label{tab:curve}
\begin{tabular}{clccc}
\toprule
\textbf{Baseline} & \textbf{Defense in place} & \textbf{Signal read} &
\textbf{Tracks executor?} & \textbf{\ais{}} \\
\midrule
B1 & Shared service account        & shared credential          & undefined & \textbf{0.0} \\
B2 & Per-agent identity            & execution-time authenticator & yes     & \textbf{1.0} \\
B3 & + RFC 8693 delegation         & delegation current actor   & no        & \textbf{0.0} \\
B4 & + tamper-evident log          & delegation current actor   & no        & \textbf{0.0} \\
\bottomrule
\end{tabular}
\end{table}

\paragraph{Why it regresses (a one-paragraph preview).}
RFC~8693's ``current actor'' (\actsub{}) is the party that \emph{requested} the
delegated authority. In a multi-agent hand-off, the agent that \emph{executes}
the action --- the one that actually presents the token at the resource --- can
differ from the one named in the token. The standard provides no field that
records the executor when it differs from the requester. RFC~8693~\S4.1's
normative \code{MUST} is scoped to the \emph{access-control decision} (``consider
only the current actor''), not to audit logging; a realistic implementation
reuses that same access-control identity for the audit record. Combined with
unbound bearer tokens (RFC~8693 inherits OAuth~2.0's default holder model, which
does not require sender-constraint) and a topology where minting precedes
execution-routing, the system logs the \emph{requester}. The standard neither
prevents nor mandates this; it simply has no place to put the executor.

\paragraph{What sets this apart.}
Unlike delegation frameworks that report detection rates (can an attack be
blocked?), AEGIS-AT treats attribution correctness itself --- does the logged
actor equal the true executor? --- as the measured dependent variable.

\paragraph{Contributions.}
\begin{enumerate}[leftmargin=1.4em,itemsep=2pt]
  \item \textbf{A metric.} The Attribution Integrity Score (\ais{}): a strict,
  three-field (\code{actor}, \code{scope}, \code{principal\_chain}) measure of
  whether a claimed audit record matches an independently observed ground truth,
  with a defect breakdown that diagnoses \emph{which} field broke (\S\ref{sec:metric}).
  \item \textbf{A benchmark.} Four defense baselines realized as configuration
  flags over a single codebase, so \ais{} differences are attributable to the
  flag and not to incidental implementation quality (\S\ref{sec:baselines}).
  \item \textbf{An independent ground-truth construction.} A harness recorder that
  observes the true executing agent out-of-band --- from the execution context,
  never from the token --- making the comparison non-circular (\S\ref{sec:recorder}).
  \item \textbf{A pre-registered, reproduced result.} The non-monotonic curve,
  predicted in the threat model before the attack code was written, reproduced
  deterministically against the real modules; the attack succeeds \emph{by
  construction} (\S\ref{sec:results}).
  \item \textbf{A structural argument and a named fix.} An analysis of why the gap
  is topology-independent, and identification of the standardized layer
  (sender-constrained tokens) hypothesized to close it, scoped as future work
  (\S\ref{sec:discussion}, \S\ref{sec:future}).
\end{enumerate}

\paragraph{Scope and honesty up front.}
This is one minimal system ($n=1$), with deterministic scripted agents (no LLM),
and Baseline~4 is attribution-only in v1 (the hash-chained tamper-evident log is
named as future work). These are deliberate boundaries that buy a clean causal
claim; we state them here and defend them in \S\ref{sec:validity} rather than
let a reviewer find them.

% ===========================================================================
\section{Background and Related Work}
\label{sec:background}
% ===========================================================================

\subsection{RFC 8693 delegation and the \texttt{act} claim}
OAuth~2.0 Token Exchange (RFC~8693) represents delegation with the \code{act}
(``actor'') claim. A token's \code{sub} names the principal whose authority is
being exercised; the \code{act} claim names who is acting on that principal's
behalf, and \code{act} can nest, forming a verifiable chain back to the original
human. RFC~8693~\S4.1 is explicit that, for access-control purposes, a consumer
\code{MUST} consider only the token's top-level claims and the party identified
as the \emph{current actor} by the (outermost) \code{act} claim; prior actors in
nested \code{act} claims are ``informational only.'' The spec's own delegation
example (\S A.2.5) and its description of the actor as ``the actor that will
wield the security token'' (\S A.2.3) implicitly assume that the party who
\emph{requested} the token and the party who \emph{wields} it at the resource are
the same entity. That assumption is the seam this paper pries open.

\subsection{The confused deputy, old and new}
The confused-deputy pattern --- a privileged component induced to misuse its
authority on behalf of a less-privileged party --- dates to
Hardy~\cite{hardy1988}. In agentic AI it has re-emerged as a high-severity
pattern: the Cloud Security Alliance's March~2026 research
note~\cite{csa2026} establishes confused-deputy prompt injection as a
high-severity threat in agent deployments and observes that when an action runs
under a trusted agent's identity, ``audit logs may look legitimate and delay
detection'' --- precisely the attribution failure we measure. The most public
production instance is \emph{Clinejection} (February~2026): a crafted GitHub
issue title drove an AI triage agent into a supply-chain compromise, an
unauthorized npm package reaching $\sim$4{,}000 developer machines in an eight-hour
window~\cite{clinejection}. The Salesloft~Drift breach
(August~2025)~\cite{salesloft} is a production demonstration of the
\emph{unbound bearer-token} weakness our Baseline~3 depends on: stolen OAuth
tokens, presented by a party that was not their legitimate holder, were used to
exfiltrate data from 700+ organizations.

\subsection{The standards call}
NIST's NCCoE concept paper \emph{Accelerating the Adoption of Software and AI
Agent Identity and Authorization} (February~2026)~\cite{nist2026} names auditing
and non-repudiation of AI agent actions as an open problem and asks how existing
identity standards (OAuth~2.0, OpenID, SPIFFE) should apply to multi-agent
delegation, including the unresolved multi-hop case. The OpenID Foundation's
response~\cite{openid2026} frames the most urgent risks as failures of
\emph{trust} --- ``Who authorised this agent to act? On whose behalf? Can that be
verified?'' --- and notes that today's deployments rely on unsigned credentials
and no clear chain of accountability. The Foundation for American Innovation's
\emph{Human-Anchored Intent-Bound Delegation} (HAID) proposal~\cite{haid}
offers a candidate fix: signed, scope-attenuating, human-anchored delegation
chains --- the kind of execution-identity binding we name as future work.

\subsection{Closest prior academic work}
\emph{The Misattribution Gap}~\cite{misattribution} measures misattribution in
agentic systems across 64 documented failures, finding that attribution systems
consistently blamed the model rather than the poisoned memory layer. It is
adjacent but a \emph{different layer}: it measures model-vs-memory
misattribution (memory poisoning, semantic norm drift), whereas AEGIS-AT
measures agent-vs-agent sibling impersonation at the delegation layer. To our
knowledge no published adversarial benchmark scores attribution integrity for
the sibling-impersonation case across a structured defense gradient; the problem
space is \emph{defensibly underexplored} rather than untouched.

\paragraph{SentinelAgent / DelegationBench v4.}
The closest adjacent framework is SentinelAgent / DelegationBench~v4~\cite{sentinelagent}: a Delegation
Chain Calculus, a non-LLM Delegation Authority Service (DAS), and a
516-scenario benchmark (150 adversarial actions across seven attack categories;
366 benign), reporting 100\% attack detection at 0\% false-positive rate with
TLA+ verification of its deterministic properties --- including \emph{forensic
reconstructibility} (P4): given any action, the hash-linked delegation chain
that authorized it is reconstructible in $O(n)$. SentinelAgent and AEGIS-AT
measure different quantities. SentinelAgent measures \emph{detection} --- does
the DAS block a policy-violating delegation? AEGIS-AT measures \emph{attribution
integrity} --- does the logged actor equal the true executor? --- on actions
that violate no policy. The two meet at one point: P4 reconstructs the lineage
of the token \emph{presented at the proxy}, which establishes who was
\emph{authorized}, not who \emph{wielded} the token at the resource, unless the
token is sender-constrained. SentinelAgent's tokens are HMAC-signed bearer
credentials with no proof-of-possession; under the re-delegation hand-off we
study (\S5), a sibling presents a token whose chain names another agent, so
hash-chain reconstruction succeeds while executor attribution fails. AEGIS-AT
therefore does not refute reconstructibility --- it \emph{bounds the
audit-attribution interpretation of P4}: a walkable chain certifies who was
authorized, not who acted. A DAS could close this by binding the token to its
holder (DPoP / mTLS) or authenticating the caller independently, but that is an
added sender-constraint, not a consequence of reconstruction. The attack also
needs none of SentinelAgent's stronger adversaries (e.g. its compromised-agent
case): it succeeds with alert-text control alone, no component misbehaving.
Among the works we compare, AEGIS-AT is the only one that treats attribution
correctness as a \emph{measured} dependent variable rather than a design
guarantee.

% ===========================================================================
\section{Threat Model}
\label{sec:threat}
% ===========================================================================

\subsection{System under test}
AEGIS-AT models a minimal SOC alert-triage pipeline (Figure~\ref{fig:arch}). A
human analyst authenticates and issues a request to a triage orchestrator. The
orchestrator delegates to two sibling subagents by minting each a scoped
RFC~8693 delegation token. Both subagents can call one shared tool, \siem{},
which records the identity of the calling agent in an action log.

\begin{itemize}[leftmargin=1.4em,itemsep=2pt]
  \item \textbf{Human principal} (\code{human:analyst}) --- authenticates once,
  originates the task, is the root of every delegation chain; does not call tools
  directly.
  \item \textbf{Orchestrator} --- a stateless RFC~8693 token-exchange endpoint. It
  validates an exchange request and mints the appropriately scoped, correctly
  nested token. It does \emph{not} read alert content to make routing decisions.
  \item \textbf{Agent-Enrich} (Subagent~A) --- the lower-consequence sibling;
  read-only context gathering under \code{siem:read}. In the attack, this is the
  innocent sibling whose identity is falsely stamped on Contain's action.
  \item \textbf{Agent-Contain} (Subagent~B) --- the higher-consequence sibling and
  the true executor; consequential response actions under \code{siem:write}.
  \item \textbf{Tool} \siem{} --- a single SOAR-style endpoint, scope-gated. The
  \code{command} parameter selects the operation; the token's \code{scope} claim
  decides whether the call is permitted. The tool reads agent identity
  \emph{only} from the verified \code{act} claim, never a self-reported field.
\end{itemize}

This is deliberately the smallest configuration in which sibling impersonation is
possible: two siblings (so impersonation has a target) and one shared tool (so
the attack has exactly one degree of freedom --- A acted, the log says B). The
privilege asymmetry between Contain and Enrich is what makes misattribution
\emph{security-relevant} rather than cosmetic.

\begin{figure}[h]
\centering
\renewcommand{\arraystretch}{1.25}
\begin{tabular}{@{}r@{\quad}l@{}}
\toprule
\multicolumn{2}{@{}l}{\textbf{Step / actor}\hfill\textbf{What happens (and at which boundary)}}\\
\midrule
\code{human:analyst}     & issues the triage request \hfill \emph{(Boundary 1)} \\
$\downarrow$             & orchestrator mints scoped RFC 8693 tokens \hfill \emph{(Boundary 2 --- attack)} \\
\code{Agent-Enrich}      & (requester, \code{siem:read}) sends a re-delegation request \\
$\downarrow$             & orchestrator mints a \code{siem:write} token \emph{naming Enrich} \\
\code{Agent-Contain}     & (true executor) wields the token, calls \siem{} \hfill \emph{(Boundary 3: all checks pass)} \\
$\downarrow$             & tool writes \code{claimed\_actor = enrich} \hfill \emph{(Boundary 4)} \\
\code{Harness recorder}  & observes \code{true\_actor = contain}, out-of-band \hfill \emph{(Boundary 5)} \\
\midrule
\textbf{AIS scorer}      & \code{claimed} $\neq$ \code{true} $\Rightarrow$ attribution defect \\
\bottomrule
\end{tabular}
\caption{The AEGIS-AT pipeline and its five trust boundaries. The attack lives at
Boundary~2 (token minting); Boundaries~3 and 4 faithfully transcribe a record
that is already wrong; Boundary~5 is the independent measurement instrument.}
\label{fig:arch}
\end{figure}

\subsection{Trust boundaries}
The system has five boundaries; naming all five lets a reviewer see that the
attack lives at exactly one of them.
\begin{description}[leftmargin=1.6em,style=nextline,itemsep=2pt]
  \item[Boundary 1 (principal $\to$ orchestrator).] The analyst authenticates via
  OIDC; their task prompt is trusted. \emph{But} the pipeline also ingests
  upstream SIEM alert content, which is definitionally untrusted --- it describes
  activity outside the perimeter and routinely carries attacker-controlled
  strings. This is the attacker's sole injection point.
  \item[Boundary 2 (orchestrator $\to$ subagent --- minting).] The attack
  boundary. Cryptographic verification is sound (the signing key is out of
  reach), but the \emph{decisional} content of the exchange request is shaped by
  upstream, attacker-influenceable input. The orchestrator honestly mints a token
  that names the wrong sibling as actor.
  \item[Boundary 3 (subagent $\to$ tool).] The tool runs full verification ---
  signature, expiry, chain integrity, scope gate, identity resolution --- and
  \emph{every check passes}. It then faithfully records the current actor from the
  verified \code{act} claim. The misattribution is not a Boundary~3 failure; it is
  a Boundary~2 outcome that Boundary~3 honestly transcribes.
  \item[Boundary 4 (tool $\to$ action log).] The claimed actor is written to the
  log. The wrong actor arrived already wrong. In Baselines~1--3 the log is a plain
  append-only file; Baseline~4's tamper-evident design (future work in v1)
  protects the integrity of the (wrong) record but cannot detect the upstream
  misattribution.
  \item[Boundary 5 (system $\to$ harness --- ground-truth recorder).] The
  measurement instrument, not a defense. It observes the true executing agent
  out-of-band, independent of any token or claim.
\end{description}

\subsection{Adversary model}
The adversary is modeled at the \emph{gray-box} level: they know the kind of
system and the kind of decisions it makes, but not its exact prompts, extraction
code, or scope strings. This matches every real-world precedent cited above.
\begin{itemize}[leftmargin=1.4em,itemsep=1pt]
  \item \textbf{Can:} control the text of one or more SIEM alerts processed by
  Enrich (hostnames, URLs, log lines, etc.); cause those alerts to enter through
  normal channels; repeat the attempt (no rate limit assumed).
  \item \textbf{Cannot:} compromise the analyst, SIEM, or upstream systems; forge
  or modify tokens (the signing key is out of reach); read or write the
  ground-truth recorder; compromise the orchestrator, subagents, or tool
  \emph{as processes} --- only the data they consume.
\end{itemize}
\textbf{Goal:} cause an action truly performed by Agent-Contain to be attributed
to Agent-Enrich, so that an analyst reviewing the log sees no record of the
high-consequence agent having taken the high-consequence action. The attack is
``covering tracks,'' not privilege escalation: the attacker hides \emph{who}
exercised authority that already legitimately existed. Critically, the attack
needs \emph{no} zero-day, no jailbreak, no instruction-override, no adversarial
embedding --- only that a genuinely containment-warranting alert flow through the
re-delegation path.

% ===========================================================================
\section{The Attribution Integrity Score}
\label{sec:metric}
% ===========================================================================

\subsection{Ground-truth and claimed schemas}
For each tool invocation the harness's ground-truth recorder writes a tuple
\[
\langle \texttt{true\_actor},\ \texttt{true\_scope},\
\texttt{true\_principal\_chain},\ \texttt{command},\ \texttt{target},\
\texttt{timestamp}\rangle,
\]
where \code{true\_actor} is the agent that actually executed the call (observed
from the executing context, not any token); \code{true\_scope} is the scope the
command genuinely requires (a static command$\to$scope map); and
\code{true\_principal\_chain} is the ordered delegation path from immediate actor
to human principal, $[\texttt{true\_actor},\ \texttt{human:analyst}]$ --- a
two-hop chain. The orchestrator does \emph{not} appear in it: it is a stateless
minting endpoint, not a delegated principal holding a token, so RFC~8693's
\code{act} claim records no hop for it.\footnote{This two-hop shape is itself a
finding in miniature. An earlier draft of the threat model recorded a three-hop
chain $[\texttt{actor},\ \texttt{orchestrator},\ \texttt{analyst}]$ --- the
natural practitioner intuition. That intuition is wrong about RFC~8693, and the
gap between how engineers reason about delegation chains and what the standard
actually records is exactly the gap that makes this misattribution surface.}

The corresponding Boundary~4 record has the parallel shape with each
\code{claimed\_*} field derived from the verified \code{act} claim of the
presented token. Records are matched between the two stores by the
$(\texttt{command}, \texttt{target}, \texttt{timestamp})$ triple.

\subsection{The metric}
For a single adversarial action $a$, define
\[
\textsf{is\_correct}(a) =
\begin{cases}
1 & \text{if } \texttt{claimed\_actor}(a) = \texttt{true\_actor}(a)\\
  & \ \wedge\ \texttt{claimed\_scope}(a) = \texttt{true\_scope}(a)\\
  & \ \wedge\ \texttt{claimed\_principal\_chain}(a) = \texttt{true\_principal\_chain}(a)\\
0 & \text{otherwise.}
\end{cases}
\]
Comparison is strict: all three fields must match. \code{principal\_chain}
comparison is ordered-list equality (a permutation, missing hop, or inserted hop
is a defect); where no delegation chain exists (opaque per-agent credentials at
Baselines~1--2), the chain is \code{None} and matches \code{None}. The
Attribution Integrity Score for a baseline $B$ is
\[
\ais(B) = \frac{1}{|A(B)|} \sum_{a \in A(B)} \textsf{is\_correct}(a),
\]
where $A(B)$ is the set of adversarial actions executed under $B$ --- tool calls
the attack actually influenced. Non-adversarial calls (setup, calibration) are
excluded from the denominator. \ais{} is reported \emph{per baseline}; the curve
across baselines is the result.

\subsection{Derived reporting}
Two diagnostics travel with each \ais{}. The \textbf{defect breakdown} records,
for each incorrect action, which of the three fields mismatched --- showing
whether a baseline breaks attribution uniformly or concentrates on one field.
The \textbf{hold rate} reports the marginal \ais{} change at each baseline
transition ($1{\to}2$, $2{\to}3$, $3{\to}4$), isolating each defense layer's
contribution.

\subsection{Why the v1 measurement is categorical}
The v1 attack is \emph{categorical, not stochastic}: under scripted deterministic
agents the misattribution succeeds by construction on every adversarial action.
The finding is therefore a curve \emph{shape}, not a frequency estimate.
Confidence intervals are degenerate by design ($\ais{} \in \{0,1\}$ per
baseline). The harness establishes this with a determinism check that proves each
baseline yields byte-identical records across repeated runs, so a single
canonical execution per baseline suffices. A probabilistic-policy sweep that
would make \ais{} a real attack-frequency estimate (with Wilson intervals) is
named as future work.

% ===========================================================================
\section{The Attack Mechanism}
\label{sec:attack}
% ===========================================================================

\subsection{The structural property exploited}
Under RFC~8693, when agent $X$ requests a delegated token to act on behalf of
principal $P$, the resulting token's \code{act} claim records $X$ as the current
actor. This is correct: the chain answers ``on whose authority, through which
delegating parties.'' But an audit log is trusted to answer a subtly different
question --- ``who \emph{performed} this action?'' In a single-agent setting the
two coincide. In a multi-agent setting where one agent requests a delegated
capability that \emph{another} agent executes, they diverge, and the divergence
is invisible to every cryptographic check because nothing was forged or
malformed.

\subsection{Token structure under Baselines 3--4}
When the orchestrator performs the re-delegation exchange, the issued token
follows RFC~8693~\S4.1 delegation semantics:
\begin{lstlisting}
sub:   "human:analyst"        <- principal (on whose behalf)
scope: "siem:write"
act: {
  sub: "agent:enrich"         <- current actor (requester / named wielder)
}
\end{lstlisting}
The chain is two hops: current actor (\code{agent:enrich}) and root principal
(\code{human:analyst}). Two properties are decisive, both following directly from
\S4.1: (i)~the current actor is the \emph{requester}, not the executor; and
(ii)~the executing agent (\code{agent:contain}) \emph{does not appear in the
token at all}. There is no field in the spec-compliant structure that records
``who wielded the token at the resource'' distinct from ``who was delegated the
authority.'' The misattribution is therefore not a property of any field being
read incorrectly; it is a property of the token having no field that can express
the divergence.

\subsection{The attack, step by step}
The attacker's only capability is to control the content of an alert Enrich
processes. The attack does not require fooling Enrich, embedding instructions, or
causing any component to misbehave --- only that a genuinely
containment-warranting alert flow through the Enrich~$\to$~re-delegation path.
\begin{enumerate}[leftmargin=1.6em,itemsep=1pt]
  \item The alert legitimately contains containment-warranting indicators (the
  attacker shaped them, but they are \emph{real} indicators). Enrich is not
  deceived; it makes the right call.
  \item Enrich correctly initiates a re-delegation request for a containment
  action, presenting its own token as the \code{actor\_token}.
  \item The orchestrator correctly mints a \code{siem:write} token whose
  \code{act} chain names Enrich (the requester).
  \item Agent-Contain --- the agent that actually holds and runs the containment
  capability --- wields the token and calls \siem{}.
  \item The tool verifies the token (every Boundary~3 check passes), reads the
  current actor (\code{enrich}), and records $\texttt{claimed\_actor} =
  \texttt{enrich}$.
  \item The ground-truth recorder observes $\texttt{true\_actor} =
  \texttt{contain}$.
  \item The scorer flags an actor mismatch: claimed Enrich, true Contain.
  Attribution has failed.
\end{enumerate}
An analyst reviewing the log sees the high-consequence containment attributed to
the read-only enrichment agent; the true executor appears nowhere in the record.

\subsection{Why it survives every objection}
\textbf{``The orchestrator has a bug.''} No --- it validated cryptographically and
built the chain from the presented \code{actor\_token}, exactly as \S4.1
specifies. \textbf{``Enrich was prompt-injected.''} No --- Enrich made the correct
decision; the result holds even if Enrich is a perfect, unfoolable agent.
\textbf{``The tool skipped a check.''} No --- every check passed; the tool
faithfully recorded the current actor. \textbf{``You didn't sanitize alert
text.''} No --- no component read identity from alert text; alert content's only
role was to legitimately trigger a containment-warranting situation. The attack
works because the delegation chain answers ``who requested'' while the audit log
is trusted to answer ``who acted,'' and in the re-delegation pattern those are
different agents.

% ===========================================================================
\section{Defense Baselines}
\label{sec:baselines}
% ===========================================================================

The benchmark measures \ais{} across four configurations, applied as config flags
over one codebase (not four implementations). Each adds one layer to the
previous; \emph{only the credential differs}. For each baseline we state the
signal the tool reads and whether it tracks the true executor.

\begin{description}[leftmargin=1.6em,style=nextline,itemsep=3pt]
  \item[Baseline 1 --- Shared service account.] All agents share one credential.
  The tool cannot distinguish callers; attribution is \emph{undefined} (not wrong
  about a specific sibling --- there is no per-agent identity to be right about).
  Predicted $\ais{} \approx 0.0$. This is the common-but-wrong status quo (shared
  API keys).
  \item[Baseline 2 --- Per-agent identity.] Each agent holds its own credential;
  attribution is bound at authentication time. Contain executes, so Contain's
  credential authenticates, so the tool records Contain. \emph{Claimed actor =
  true actor.} Predicted $\ais{} \approx 1.0$. Attribution is correct not by
  design sophistication but because here the \emph{executor is the
  authenticator}.
  \item[Baseline 3 --- + RFC 8693 delegation.] The tool now resolves the claimed
  actor from the delegation chain's current actor (\actsub{}~$=$~Enrich, the
  requester), as \S4.1 mandates for the access-control decision. Under unbound
  bearer tokens, Contain lifts the token minted for Enrich and presents it, and
  nothing in the protocol detects the substitution. Predicted $\ais{} \approx
  0.0$. This is the central result: signed delegation \emph{regresses}
  attribution relative to per-agent identity, by following the standard
  correctly.
  \item[Baseline 4 --- + tamper-evident log.] Same signal as B3. Tamper-evidence
  protects the integrity of the recorded entry but does not change \emph{what} is
  recorded; the wrong actor was committed upstream at mint time. Predicted
  $\ais{} \approx 0.0$, unchanged from B3. (In v1 this module is not built ---
  B4's attribution equals B3 by construction; a real hash-chained log, which
  tests log \emph{integrity}, a separate metric, is future work.)
\end{description}

The two ingredients of the regression are neither sufficient alone: (i)~unbound
bearer tokens carrying a current-actor claim (RFC~8693 inherits OAuth~2.0's
default holder model and does not require sender-constraint), and (ii)~a
multi-agent hand-off where the wielder differs from the issuer-named actor (the
orchestrator must name the requester because the executor is undetermined at mint
time). \S4.1's \code{MUST} is scoped to access control and is silent on audit; a
realistic implementation reuses the access-control identity for the audit record,
and that identity is necessarily the named requester.

% ===========================================================================
\section{Implementation}
\label{sec:impl}
% ===========================================================================

The system is implemented in Python ($\sim$6 core modules), using PyJWT for
RS256-signed tokens and the \code{cryptography} library for key generation. The
agents are scripted (deterministic) by design --- this isolates the
delegation-layer failure from model behavior.

\begin{description}[leftmargin=1.6em,style=nextline,itemsep=2pt]
  \item[\code{auth/tokens.py}] Mints and verifies RFC~8693 tokens by hand:
  \code{mint\_initial\_token}, \code{exchange\_token} (nests \code{act}, narrows
  scope, refuses to widen), \code{verify\_token}, and \code{actor\_chain}, which
  returns the path \emph{current-actor-first, root-principal-last}. Identity
  resolution reads \code{chain[0]} --- the current actor (\S4.1), never the
  innermost subject (the root principal).
  \item[\code{policy/scope\_map.py}] The shared command$\to$scope contract
  (\code{keyword\_search}~$\to$~\code{siem:read}, \code{isolate\_host}~$\to$~%
  \code{siem:write}). Imported by both the tool and the recorder so they cannot
  drift. Unknown commands raise --- a fail-loud policy error, not a default.
  \item[\code{tools/siem\_action.py}] The scope-gated tool (Boundary~3): signature
  + expiry (PyJWT), chain integrity, scope gate, identity resolution. For opaque
  per-agent credentials (Baselines~1--2) it reads the authenticating agent and
  sets \code{principal\_chain}~$=$~\code{None}; for JWTs (Baselines~3--4) it takes
  the \code{act}-claim path. Discrimination is on credential \emph{structure}
  only.
  \item[\code{orchestrator/orchestrator.py}] A thin RFC~8693 validator/minter
  (Boundary~2). The only validation it adds beyond \code{exchange\_token} is that
  the \code{sub} must name a human principal. It does not route on or read alert
  content; it does not appear in the minted chain.
  \item[\code{harness/recorder.py}] The independent ground-truth recorder
  (Boundary~5); see \S\ref{sec:recorder}.
  \item[\code{harness/scorer.py}] The \ais{} metric, defect breakdown, and Wilson
  CI; plus \code{is\_non\_monotonic(curve)}, the named predicate that pins the
  \S\ref{sec:baselines} headline claim in one place.
  \item[\code{harness/sweep.py}] The baseline switch: \code{run(baseline)} mints
  the per-baseline credential (the only thing that varies), runs the canonical
  attack once on a thread named for the true executor, and scores it;
  \code{verify\_deterministic} and \code{emit\_curve} compose it.
\end{description}

\subsection{The independent ground-truth recorder}
\label{sec:recorder}
The recorder wraps \siem{} and observes the calling thread's name as
\code{true\_actor} \emph{before} forwarding the call. Its independence rests on
three axes: a \textbf{process boundary} (it runs in the harness, distinct from
the agent contexts the attacker influences); \textbf{credential isolation} (no
in-system component holds ground-truth write credentials); and \textbf{causal
precedence} (it records before the tool's verification logic runs). The
attacker's only capability --- alert-content control --- cannot cross any of
these. If this independence failed, the measurement would be circular; it is
established by construction. A v1 honesty note: the recorder uses Python's
\code{threading.current\_thread().name} as a proxy for OS process identity; a
true process boundary via \code{multiprocessing}/\code{os.getpid()} is v2
hardening. Within the threat model (the adversary controls alert text only, not
agent code), the proxy holds.

% ===========================================================================
\section{Experimental Methodology}
\label{sec:method}
% ===========================================================================

\paragraph{Pre-registration.} Every \ais{} value (B1~$\approx$~0, B2~$\approx$~1.0,
B3~$\approx$~0, B4~$\approx$~0) was committed in the threat model \emph{before}
the attack code was written. The test suite asserts the measured curve against
these predictions; a contradicted prediction is to be reported as a finding, not
reconciled. This is the structural difference between a benchmark and a demo.

\paragraph{Determinism as a checked property.}
\code{verify\_deterministic(baseline, k)} runs each baseline $k=5$ times under a
fixed clock and asserts byte-identical \code{(claimed, truth)} records, raising
on the first divergent field. \code{emit\_curve} gates on this check before
emitting any \ais{} value, so the curve is only produced once determinism is
proven rather than asserted. A fault-injection test feeds a deliberately drifting
clock to confirm the check actually catches divergence.

\paragraph{Comparability by construction.}
Because the four baselines are config flags over one codebase --- identical tool,
recorder, and scorer code, differing only in the credential --- an \ais{}
difference is attributable to the credential flag, not to incidental
implementation quality. This is what makes the four numbers a curve rather than
four anecdotes.

\paragraph{Reproduction.}
\begin{lstlisting}
pip install -r aegis-at/requirements.txt
pytest tests/core -v          # 59 tests
bash scripts/check.sh         # lint, format, invariant greps, test gate
\end{lstlisting}
The full pipeline reproduces in seconds.

% ===========================================================================
\section{Results}
\label{sec:results}
% ===========================================================================

\subsection{The curve}
The measured curve reproduces the pre-registered prediction exactly
(Table~\ref{tab:results}). \code{is\_non\_monotonic(curve)} returns \code{True}:
$B2 > B1$, $B2 > B3$, and $B4 = B3$ (byte-identical, since B4 runs B3's code
path in v1).

\begin{table}[h]
\centering
\caption{Measured \ais{} with defect breakdown. Numerator/denominator are over
the single canonical adversarial action per baseline. The defect breakdown names
the claimed vs.\ true actor and which fields mismatched.}
\label{tab:results}
\begin{tabular}{clcccl}
\toprule
\textbf{B} & \textbf{Defense} & \textbf{\ais{}} & \textbf{num/den} &
\textbf{Predicted} & \textbf{Defect (claimed $\to$ true)} \\
\midrule
B1 & Shared account     & 0.0 & 0/1 & $\approx 0$   & actor: \code{svc:soar} $\to$ \code{agent:contain} \\
B2 & Per-agent identity & 1.0 & 1/1 & $\approx 1.0$ & --- (perfect) \\
B3 & + RFC 8693 deleg.  & 0.0 & 0/1 & $\approx 0$   & actor + chain: \code{agent:enrich} $\to$ \code{agent:contain} \\
B4 & + tamper-evident   & 0.0 & 0/1 & $\approx 0$   & actor + chain: \code{agent:enrich} $\to$ \code{agent:contain} \\
\bottomrule
\end{tabular}
\end{table}

\subsection{Reading the defect breakdown}
The defect breakdown is the cleanest evidence of the mechanism. At B3 and B4 the
system claims \code{agent:enrich} (the requester named in \actsub{}) while the
true executor is \code{agent:contain}. The defect flags the \code{actor}
\emph{and} \code{principal\_chain} fields together --- never \code{scope}. This
confirms a property predicted in the threat model: the actor and principal-chain
defects are \emph{correlated}, because both flag precisely when Enrich occupies
the current-actor position. The correlation is a true property of the attack, not
a metric artifact. At B1 the single defect is on \code{actor} alone
(\code{svc:soar} vs.\ \code{contain}), with \code{principal\_chain}~$=$~\code{None}
on both sides (no chain pre-delegation) --- the ``undefined attribution'' case.

\subsection{Test and gate status}
The full suite is \textbf{59 tests, all passing}; the mechanical gate
(\code{scripts/check.sh}) exits~0, enforcing the greppable invariants (the tool
is named \siem{} everywhere; identity resolves to the most-recent actor, never
the ``innermost'' subject), lint (ruff), and format (black). The auth smoke test
demonstrates chain nesting, scope narrowing, and forgery rejection. B4 equals B3
byte-for-byte, confirming that tamper-evidence is orthogonal to attribution: the
wrong actor is committed before any logging layer sees the entry.

% ===========================================================================
\section{Discussion}
\label{sec:discussion}
% ===========================================================================

\paragraph{This is not a logic bug a careful engineer would fix.}
The honest defense is topological, not ``the spec forces it.'' The fix a reviewer
reaches for is ``name the wielder, not the requester.'' But in the realistic
topology, the orchestrator mints the delegated token and hands it downstream,
where routing decides which sibling executes; the wielder is \emph{not determined
at mint time}, so the orchestrator cannot place it in the current-actor claim.
Naming the requester is what any orchestrator must do when minting precedes
execution-routing. Closing the gap requires binding identity at execution time
--- the wielder re-exchanges the token on receipt, naming itself current actor,
or sender-constrained tokens prevent Contain from presenting Enrich's token at
all. Both are exactly the execution-identity binding we name as future work
(\S\ref{sec:future}). The objection does not dismiss the result; it identifies
the missing layer the result measures the absence of.

\paragraph{The gap is structural, not merely adversarial.}
The misattribution is a \emph{latent property of the re-delegation pattern}: any
containment action re-delegated through Enrich produces the same wrong-actor
record, attack or no attack. The attacker does not create the vulnerability; the
attacker makes a pre-existing weakness profitable by triggering it deliberately
and repeatedly. v1 scopes the \ais{} denominator to attacker-triggered actions
for a clean measurement; an expanded denominator over all re-delegated
containment actions (v2) would measure the structural property more faithfully.

\paragraph{Delegation's auditability benefit is conditional.}
This result \emph{qualifies} rather than rejects the emerging industry
narrative that delegation beats impersonation for accountability~\cite{redhat2026, okta2026}.
Delegation does beat shared service accounts (B1) on lineage and least
privilege; what AEGIS-AT shows is that delegation's \emph{audit-attribution}
advantage is conditional --- absent sender-constraint, B3 attributes worse than
the simpler per-agent identity of B2, even while every delegation check passes.
We are not arguing impersonation is preferable; we are bounding when delegation
actually delivers the accountability it is promoted for.

\paragraph{Tamper-evident logging protects a wrong answer.}
Baseline~4 underperforms a naive expectation, and this is the point: the second
primitive standards bodies emphasize for non-repudiation cannot recover correct
attribution, because the misattribution is established \emph{before} logging. It
cryptographically preserves the wrong record. Stating this before the result
forecloses the reading ``your tamper-evident log is broken'' --- it is working
correctly; it is simply not the right defense for this attack class. That
negative result is part of the contribution.

% ===========================================================================
\section{Limitations and Validity Threats}
\label{sec:validity}
% ===========================================================================

We list the ways the result could be wrong or unconvincing and the mitigation for
each. The first is a genuine limitation we concede rather than defeat.

\begin{enumerate}[leftmargin=1.6em,itemsep=3pt]
  \item \textbf{Generalization ($n=1$) --- conceded.} The curve is demonstrated on
  one minimal pipeline with one orchestrator design and one re-delegation
  topology. v1 establishes the gap \emph{exists}, is triggerable by a realistic
  gray-box adversary, and survives the two emphasized defenses --- \emph{in this
  system}. What makes it meaningful despite $n=1$ is \emph{why} it appears: it
  follows from \S4.1 scoping its \code{MUST} to access control (reused for audit)
  plus the standard's implicit assumption that the current actor and executor
  coincide. That assumption is topology-independent; wherever a multi-agent system
  separates the requester of a capability from its executor, the gap should
  appear. Establishing prevalence across real architectures is the primary item
  of future work.
  \item \textbf{``The system is a toy.''} Minimal by design, but every structural
  element maps to a documented real-world pattern: alert-as-injection-vector
  (Clinejection, Log4Shell, Splunk~XSS), data-driven SOAR routing, and RFC~8693
  delegation as the recommended non-repudiation primitive. Minimality isolates the
  one degree of freedom (requester~$\neq$~executor) without confounds.
  \item \textbf{``Ground truth isn't independent of the log.''} Independence is
  established by construction along three axes (\S\ref{sec:recorder}); the
  attacker's only capability cannot cross any of them. Honest sub-limitation:
  established by argument, not formal verification.
  \item \textbf{``Baselines aren't a fair comparison.''} They are config flags over
  one codebase; differences cannot be attributed to incidental quality. The most
  attackable point --- B2 reaching $\approx$~1.0 --- rests on a stated model
  (attribution binds to the execution-time authenticator) with the
  authorization-vs-attribution distinction made explicit.
  \item \textbf{``Scripted, not real LLM behavior.''} Deliberate: the attack is
  designed to be independent of agent reasoning quality (Enrich's escalation is
  the \emph{correct} response to a genuinely warranting alert), removing
  model-specific confounds. It also means v1 does not measure interaction with
  imperfect agent decisions --- named out of scope.
  \item \textbf{Thread proxy for process identity.} The recorder uses
  \code{threading.current\_thread().name} as a v1 proxy; a renamed thread could
  spoof ground truth, excluded by the adversary model (alert text only). v2 uses
  \code{multiprocessing}/\code{os.getpid()}.
  \item \textbf{Baseline 4 is attribution-only in v1.} B4's attribution equals B3
  by construction; a real hash-chained tamper-evident log (testing log
  \emph{integrity}, a separate metric) is future work.
\end{enumerate}

% ===========================================================================
\section{Future Work}
\label{sec:future}
% ===========================================================================

\begin{itemize}[leftmargin=1.4em,itemsep=2pt]
  \item \textbf{Baseline 5 --- sender-constrained tokens.} The gap requires unbound
  bearer tokens. DPoP (RFC~9449) and mutual-TLS-bound tokens (RFC~8705) bind a
  token to its holder; under either, Contain cannot present Enrich's token and
  must obtain its own, so the current actor would track the executor. Whether this
  recovers the curve is the primary defensive item of future work --- effectively
  a Baseline~5. Agentic JWT~\cite{agenticjwt} is one concrete design in this
  direction --- per-agent proof-of-possession keys plus intent/delegation claims
  binding an API call to a registered agent and workflow step. In AEGIS-AT terms
  it is a plausible Baseline~5 instantiation: it proposes the sender-constraint
  but does not measure the attribution-integrity regression that motivates it.
  Pairing such a protocol with the \ais{} metric --- \emph{does sender-constraint
  recover the curve at B5?} --- is the natural next experiment.
  \item \textbf{Prevalence across topologies.} Reproduce the mechanism in
  additional multi-agent architectures to move from ``shown in one instance'' to
  ``prevalent across deployments.''
  \item \textbf{Stochastic policy / expanded denominator.} A probabilistic agent
  policy under which \ais{} becomes a real attack-frequency estimate (Wilson
  intervals, $N \geq 100$ per baseline), and a denominator over all re-delegated
  containment actions to measure the latent structural property directly.
  \item \textbf{Tamper-evident log integrity.} Implement the hash-chained, signed
  log to measure log \emph{integrity} as a metric distinct from attribution.
  \item \textbf{Other sibling-impersonation variants.} Delegation forgery,
  scope-attenuation bypass, audit-log tampering, principal laundering --- each a
  distinct mode, backlogged for v2.
\end{itemize}

% ===========================================================================
\section{Conclusion}
\label{sec:conclusion}
% ===========================================================================

Adding the industry-standard delegation mechanism to a correctly-functioning
multi-agent system can make audit attribution \emph{worse}, not better. AEGIS-AT
measures this as a non-monotonic Attribution Integrity curve: perfect under
per-agent identity, zero once RFC~8693 delegation is added, still zero with
tamper-evident logging. The cause is structural --- RFC~8693's current actor names
the requester, the executor can differ, and the standard has no field for the
executor --- so the result is not a bug to be patched but a missing layer to be
added: execution-identity binding via sender-constrained tokens. The benchmark is
small, deterministic, pre-registered, and reproducible; its value is in the
trustworthiness of the measurement, not the breadth of coverage. As standards
bodies move to make RFC~8693 the backbone of AI agent non-repudiation, AEGIS-AT
is a concrete, falsifiable warning that delegation alone does not buy
accountability in the multi-agent case --- and a precise statement of what does.

% ===========================================================================
\section*{Reproducibility and Artifact Availability}
\addcontentsline{toc}{section}{Reproducibility and Artifact Availability}
% ===========================================================================
The complete benchmark --- threat model, six core modules, harness, and 59 tests
--- is released. Code (\code{aegis-at/}, \code{tests/}, \code{scripts/}) is
licensed Apache-2.0; documentation (\code{Documents/}) is licensed CC~BY~4.0.
Every \ais{} value is asserted in the test suite against the curve predicted in
the threat model before the attack code was written. The pipeline reproduces in
seconds via \code{pytest tests/core} and \code{bash scripts/check.sh}.

% ===========================================================================
\begin{thebibliography}{99}
\addcontentsline{toc}{section}{References}

\bibitem{rfc8693} M.~Jones, A.~Nadalin, B.~Campbell, J.~Bradley, and C.~Mortimore,
\emph{OAuth 2.0 Token Exchange}, RFC~8693, IETF, January 2020.
\url{https://www.rfc-editor.org/rfc/rfc8693}.

\bibitem{rfc9449} D.~Fett, B.~Campbell, J.~Bradley, T.~Lodderstedt, M.~Jones, and
D.~Waite, \emph{OAuth 2.0 Demonstrating Proof of Possession (DPoP)}, RFC~9449,
IETF, September 2023. \url{https://www.rfc-editor.org/rfc/rfc9449}.

\bibitem{rfc8705} B.~Campbell, J.~Bradley, N.~Sakimura, and T.~Lodderstedt,
\emph{OAuth 2.0 Mutual-TLS Client Authentication and Certificate-Bound Access
Tokens}, RFC~8705, IETF, February 2020.
\url{https://www.rfc-editor.org/rfc/rfc8705}.

\bibitem{nist2026} H.~Booth, W.~Fisher, R.~Galluzzo, and J.~Roberts,
\emph{Accelerating the Adoption of Software and Artificial Intelligence Agent
Identity and Authorization}, NIST NCCoE Concept Paper (Initial Public Draft),
February~5, 2026.
\url{https://www.nccoe.nist.gov/publications/other/accelerating-adoption-software-and-ai-agent-identity-and-authorization-concept}.

\bibitem{openid2026} OpenID Foundation, \emph{OIDF Responds to NIST on AI Agent
Security}, March~11, 2026.
\url{https://openid.net/oidf-responds-to-nist-on-ai-agent-security/}.

\bibitem{csa2026} Cloud Security Alliance AI Safety Initiative, \emph{Confused
Deputy Attacks on Autonomous AI Agents} (Research Note), March~23, 2026.
\url{https://labs.cloudsecurityalliance.org/research/csa-research-note-ai-agent-confused-deputy-prompt-injection/}.

\bibitem{haid} Foundation for American Innovation, \emph{Human-Anchored
Intent-Bound Delegation (HAID) for AI Agents}, submitted to NIST, April~2, 2026.
\url{https://www.thefai.org/posts/human-anchored-intent-bound-delegation-for-ai-agents}.

\bibitem{clinejection} Snyk, \emph{How ``Clinejection'' Turned an AI Bot into a
Supply Chain Attack}, February 2026.
\url{https://snyk.io/blog/cline-supply-chain-attack-prompt-injection-github-actions/}.

\bibitem{salesloft} Google Threat Intelligence Group, \emph{Data Theft from
Salesforce Instances via the Salesloft Drift Integration (UNC6395)}, August 2025.
\url{https://cloud.google.com/blog/topics/threat-intelligence/data-theft-salesforce-instances-via-salesloft-drift}.

\bibitem{misattribution} T.~Ahad, I.~Hossain, M.~J.~Alam, S.~Puppala,
S.~B.~Alam, and S.~Talukder, \emph{The Misattribution Gap: When Memory Poisoning
Looks Like Model Failure in Agentic AI Systems}, arXiv:2605.22842, May 2026.
\url{https://arxiv.org/abs/2605.22842}.

\bibitem{hardy1988} N.~Hardy, \emph{The Confused Deputy (or why capabilities might
have been invented)}, ACM SIGOPS Operating Systems Review, 22(4), 1988.

\bibitem{sentinelagent} K.~S.~R.~Patil, \emph{SentinelAgent: Intent-Verified Delegation Chains for
Securing Federal Multi-Agent AI Systems}, arXiv:2604.02767, April 2026.
\url{https://arxiv.org/abs/2604.02767}

\bibitem{agenticjwt} A.~Goswami, \emph{Agentic JWT: Secure Delegation for Agentic AI},
arXiv:2509.13597, September 2025. \url{https://arxiv.org/abs/2509.13597}

\bibitem{redhat2026} Red Hat, \emph{Zero Trust for AI Agents: Why Delegation Beats Impersonation},
Red Hat Emerging Tech, May~21, 2026.
\url{https://next.redhat.com/2026/05/21/zero-trust-for-ai-agents-why-delegation-beats-impersonation/}

\bibitem{okta2026} Okta, \emph{Agent Security: Securing the Delegation Chain}, Okta AI Security Blog, 2026.
\url{https://www.okta.com/blog/ai/agent-security-delegation-chain/}

\end{thebibliography}

\end{document}
